\documentclass[12pt,a4paper]{article}

% ------------------------------------------------
% Pacotes
% ------------------------------------------------
\usepackage[utf8]{inputenc}
\usepackage[T1]{fontenc}
\usepackage[brazil]{babel}
\usepackage{graphicx}
\usepackage[table]{xcolor}
\usepackage{array}
\usepackage{tabularx}
\usepackage{geometry}
\usepackage{fancyhdr}
\usepackage{mathptmx} % Times New Roman
\usepackage{enumitem}\usepackage{amsfonts}
\usepackage{amsmath} % Recomendado para o comando \text{}
\usepackage{amssymb}
\usepackage{tikz}
\usetikzlibrary{positioning, shapes.geometric}

% TikZ Styles for ER Diagrams
\tikzset{
    block/.style={rectangle, draw, thick, minimum width=2cm, minimum height=1cm, fill=white},
    diamond_node/.style={diamond, draw, thick, minimum width=2cm, minimum height=1cm, fill=white, aspect=2},
    attribute/.style={ellipse, draw, thin, minimum width=1.5cm, minimum height=0.6cm, fill=white, font=\scriptsize},
    line/.style={draw, thick},
}
% ------------------------------------------------ 
% Configuração de página
% ------------------------------------------------
\geometry{
    a4paper,
    left=2cm,
    right=2cm,
    top=2cm,
    bottom=2cm,
    headheight=15pt % Adicione isso
}% Caixa para resposta discursiva
\newcommand{\resposta}[1]{
\noindent
\framebox[\textwidth]{
\begin{minipage}{0.98\textwidth}
\vspace{#1}
\end{minipage}
}
}

\pagestyle{fancy}
\fancyhf{}
\renewcommand{\headrulewidth}{0pt}

\setlength{\parindent}{0pt}

% ------------------------------------------------
% Comandos auxiliares
% ------------------------------------------------
\newcolumntype{L}[1]{>{\raggedright\arraybackslash}p{#1}}
\newcolumntype{C}[1]{>{\centering\arraybackslash}p{#1}}

% ------------------------------------------------
% Documento
% ------------------------------------------------
\begin{document}

% ------------------------------------------------
% Cabeçalho - Corrigido para evitar aninhamento problemático
% ------------------------------------------------
\begin{tabularx}{\textwidth}{@{} L{4cm} X @{}}
    \raisebox{-0.5cm}{\includegraphics[width=4cm]{Universidade_de_Rio_Verde.jpg}} &
    \hfill
    \begin{tabular}{|>{\centering\arraybackslash}m{10cm}|>{\raggedright\arraybackslash}p{1.8cm}|}
        \hline
        \cellcolor{black}\color{white}\rule{0pt}{1.2cm}\textbf{AVALIAÇÃO 3 - GABARITO} 
        & 
        \scriptsize Nota \tabularnewline
        \hline
    \end{tabular} 
\end{tabularx}
\vspace{0.3cm}

% ------------------------------------------------  
% Tabela principal - Estrutura limpa
% ------------------------------------------------
\begin{tabularx}{\textwidth}{|X|L{2cm}|}
\hline
\textbf{Curso} & \textbf{Campus} \\
\textit{Engenharia de Software} & \textit{Rio Verde} \\
\hline
\multicolumn{2}{|l|}{\textbf{Disciplina}} \\  
\multicolumn{2}{|l|}{ \textit{Banco de dados}} \\
\hline
\textbf{Nome do(a) acadêmico(a)} & \textbf{Valor} \\
& 3,0 \\
\hline
\textbf{Professor} & \textbf{Turma} \\
\textit{ Sergio Souza Novak} & \textit{B} \\
\hline
\textbf{Data da Avaliação} & \textbf{Código} \\
\textit{23/06/2026} & \textit{ESW417} \\
\hline
\end{tabularx}
 
% ------------------------------------------------
% Barra de Atenção
% ------------------------------------------------
\vspace{0.2cm}

\begin{tabularx}{\textwidth}{|X|}
\hline
\cellcolor{black}\color{white}
\centering \textbf{ATENÇÃO:} Somente serão passíveis de \textbf{REVISÃO} avaliações resolvidas à \textbf{TINTA}. \tabularnewline
\hline
\end{tabularx}

\vspace{0.3cm}

\textbf{ORIENTAÇÕES PARA A RESOLUÇÃO}  

\begin{itemize}
    \item A avaliação é individual, sem consulta;
    \item A interpretação de texto faz parte da avaliação;
\end{itemize}

\vspace{0.5cm}

\begin{center}
    \rule{\textwidth}{0.4pt}
    \textbf{QUESTÕES}
\end{center}

\vspace{0.2cm}

\begin{enumerate}
    \item  Em muitas aplicações web e mobile, a exibição de grandes volumes de dados é feita de forma dividida através de páginas (paginação) para melhorar o desempenho e a usabilidade. Considere um banco de dados de uma loja virtual com o seguinte modelo de entidade para os \textbf{produtos}:
    
    \vspace{0.3cm}
    \begin{center}
    \begin{tikzpicture}
        % Entidade
        \node[block] (produtos) at (0,0) {\textbf{produtos}};
        
        % Atributos
        \node[attribute] (id) at (-2.5, 1.2) {\underline{id}};
        \node[attribute] (nome) at (2.5, 1.2) {nome};
        \node[attribute] (preco) at (-2.5, -1.2) {preco};
        \node[attribute] (categoria) at (2.5, -1.2) {categoria};
        
        % Ligações
        \draw[line] (produtos) -- (id);
        \draw[line] (produtos) -- (nome);
        \draw[line] (produtos) -- (preco);
        \draw[line] (produtos) -- (categoria);
    \end{tikzpicture}
    \end{center}
    \vspace{0.2cm}
    
    Escreva uma consulta SQL que retorne o \texttt{nome} e o \texttt{preco} dos produtos para exibir a \textbf{3ª página} de resultados, assumindo que cada página mostra exatamente \textbf{15 produtos}. Os resultados devem estar ordenados pelo \texttt{preco} de forma decrescente.
    
    \vspace{0.2cm}
    \noindent
    \framebox[\textwidth]{
    \begin{minipage}{0.98\textwidth}
    \vspace{0.2cm}
    \textcolor{blue}{\textbf{Gabarito:}}
    
    \textcolor{blue}{\texttt{SELECT nome, preco FROM produtos ORDER BY preco DESC LIMIT 15 OFFSET 30;}}
    
    \textcolor{blue}{(Também aceito o uso de \texttt{LIMIT 30, 15})}
    \vspace{0.2cm}
    \end{minipage}
    }

    \vspace{0.5cm}
    


\item  \textbf{Leia o texto de apoio abaixo para responder à questão:}

\begin{center}
\fbox{\parbox{0.95\textwidth}{
    \small
    \textbf{Ações Referenciais em SQL} \\
    Ações referenciais são regras aplicadas às restrições de chave estrangeira (FOREIGN KEY). Elas determinam como o banco de dados deve se comportar quando um registro "pai" (referenciado) é excluído ou atualizado. Servem para manter a integridade referencial de forma automática. As 5 ações padrão disponíveis no SQL são:
    \vspace{-0.2cm}
    \begin{enumerate}[leftmargin=*, label=\arabic*.]
        \item \textbf{CASCADE}: Propaga a alteração. Se o registro pai for deletado ou atualizado, todos os registros filhos correspondentes são automaticamente deletados ou atualizados.
        \item \textbf{RESTRICT}: Bloqueia a operação. Impede que o registro pai seja deletado ou alterado caso exista algum registro filho que dependa dele.
        \item \textbf{NO ACTION}: Semelhante ao RESTRICT na maioria dos SGBDs, rejeitando a alteração quando existem registros dependentes.
        \item \textbf{SET NULL}: Define a chave estrangeira dos registros filhos como NULL. Para isso, a coluna não pode possuir a restrição NOT NULL.
        \item \textbf{SET DEFAULT}: Define a chave estrangeira dos registros filhos para um valor padrão previamente estabelecido.
    \end{enumerate}
}}
\end{center}

Considere o diagrama lógico e os dados das tabelas \texttt{Departamento} e \texttt{Funcionario} abaixo, em que cada funcionário pertence a um departamento por meio da chave estrangeira \texttt{idDepartamento}:

\begin{center}
\begin{tikzpicture}[
    table/.style={rectangle, draw, thick, align=left, text width=3.5cm, inner sep=5pt},
    header/.style={rectangle, draw, thick, align=center, text width=3.5cm, fill=gray!20, font=\bfseries, inner sep=5pt}
]
    % Departamento
    \node[header] (dep_h) at (0,0) {Departamento};
    \node[table, below=0pt of dep_h] (dep_b) {
        \underline{idDepartamento} \\
        nome
    };
    
    % Funcionario
    \node[header] (func_h) at (6,0) {Funcionario};
    \node[table, below=0pt of func_h] (func_b) {
        \underline{idFuncionario} \\
        nome \\
        idDepartamento (FK)
    };
    
    % Relacionamento
    \draw[-stealth, thick] (func_b.west) -- node[above] {\scriptsize \textit{referencia}} (dep_b.east);
\end{tikzpicture}
\end{center}

\begin{center}
\begin{tabular}{|c|c|}
\hline
\textbf{Departamento} & \\
\hline
idDepartamento & nome \\
\hline
1 & Financeiro \\
2 & Recursos Humanos \\
3 & TI \\
\hline
\end{tabular}
\hspace{1cm}
\begin{tabular}{|c|c|c|}
\hline
\textbf{Funcionario} & & \\
\hline
idFuncionario & nome & idDepartamento \\
\hline
101 & Ana & 1 \\
102 & Bruno & 3 \\
103 & Carla & 3 \\
104 & Daniel & 2 \\
\hline
\end{tabular}
\end{center}

Suponha que o departamento de código \texttt{3} (TI) seja removido da tabela \texttt{Departamento}. Relacione cada ação referencial da coluna da esquerda com o comportamento correspondente na coluna da direita.

\begin{center}
\begin{tabular}{|c|p{8cm}|}
\hline
\textbf{Ação Referencial} & \textbf{Comportamento após remover o departamento 3} \\
\hline
(1) CASCADE & ( \textcolor{blue}{\textbf{2}} ) A exclusão do departamento é impedida porque existem funcionários associados. \\
\hline
(2) RESTRICT & ( \textcolor{blue}{\textbf{4}} ) Os funcionários Bruno e Carla permanecem cadastrados, mas seu campo \texttt{idDepartamento} passa a valer \texttt{NULL}. \\
\hline
(3) NO ACTION & ( \textcolor{blue}{\textbf{1}} ) Os funcionários Bruno e Carla são removidos automaticamente. \\
\hline
(4) SET NULL & ( \textcolor{blue}{\textbf{3}} ) A exclusão é rejeitada devido à existência de registros dependentes. \\
\hline
(5) SET DEFAULT & ( \textcolor{blue}{\textbf{5}} ) O campo \texttt{idDepartamento} dos funcionários Bruno e Carla assume um valor padrão previamente definido. \\
\hline
\end{tabular}
\end{center}

Assinale a sequência correta de preenchimento dos parênteses.

\begin{enumerate}[label=\alph*)]
    \item 2 -- 4 -- 3 -- 1 -- 5
    \item 1 -- 4 -- 2 -- 3 -- 5
    \item[\textbf{\textcolor{blue}{c)}}] \textcolor{blue}{\textbf{2 -- 4 -- 1 -- 3 -- 5}}
\item 2 -- 5 -- 1 -- 3 -- 4
\item 3 -- 4 -- 1 -- 2 -- 5
\end{enumerate}

    \vspace{0.5cm}

    \item  Considere o seguinte modelo de entidade das tabelas \texttt{funcionarios} e \texttt{log\_sistema}:
    
    \vspace{0.3cm}
    \begin{center}
    \begin{tikzpicture}
        % Entidade funcionarios
        \node[block] (func) at (0,0) {\textbf{funcionarios}};
        \node[attribute] (id_func) at (-2.5, 1.2) {\underline{id}};
        \node[attribute] (nome_func) at (2.5, 1.2) {nome};
        \node[attribute] (cargo) at (-2.5, -1.2) {cargo};
        \node[attribute] (salario) at (2.5, -1.2) {salario};
        
        \draw[line] (func) -- (id_func);
        \draw[line] (func) -- (nome_func);
        \draw[line] (func) -- (cargo);
        \draw[line] (func) -- (salario);
        
        % Entidade log_sistema
        \node[block] (log) at (8,0) {\textbf{log\_sistema}};
        \node[attribute] (id_log) at (5.5, 1.2) {\underline{idLog}};
        \node[attribute] (acao) at (10.5, 1.2) {acao};
        \node[attribute] (data) at (5.5, -1.2) {dataHora};
        \node[attribute] (tabela) at (10.5, -1.2) {tabelaAlvo};
        
        \draw[line] (log) -- (id_log);
        \draw[line] (log) -- (acao);
        \draw[line] (log) -- (data);
        \draw[line] (log) -- (tabela);
    \end{tikzpicture}
    \end{center}
    \vspace{0.3cm}

    Existe a necessidade de popular automaticamente a tabela de auditoria (\texttt{log\_sistema}) toda vez que um registro for \textbf{inserido} ou \textbf{apagado} da tabela \texttt{funcionarios}. Assinale a alternativa que apresenta a configuração estrutural e sintática \textbf{CORRETA} de um \textbf{Trigger} (Gatilho) para essa finalidade:
    
    \begin{enumerate}[label=\alph*)]
        \item O gatilho deve ser criado com \texttt{BEFORE SELECT ON funcionarios}, pois a leitura dos dados deve obrigatoriamente anteceder a gravação do log no sistema.
        \item[\textbf{\textcolor{blue}{b)}}] \textcolor{blue}{\textbf{O gatilho deve utilizar a estrutura \texttt{AFTER INSERT OR DELETE ON funcionarios FOR EACH ROW}, executando internamente um \texttt{INSERT} na tabela de log para cada linha afetada.}}
        \item A cláusula obrigatória para tabelas de auditoria é \texttt{FOR ALL TABLES}, garantindo que o gatilho seja acionado independentemente da tabela que sofreu a ação.
        \item O código de um gatilho não pode executar comandos \texttt{INSERT INTO log\_sistema}, pois gatilhos são objetos restritos exclusivamente à validação de dados em memória.
        \item Para garantir que o log seja populado, o desenvolvedor deve disparar o gatilho manualmente através da aplicação usando o comando \texttt{EXECUTE TRIGGER trg\_log}.
    \end{enumerate}

    \vspace{0.5cm}

    \item  Os índices são estruturas fundamentais para otimizar o tempo de resposta das consultas. Sobre o índice \textbf{B-Tree} (Árvore B) e como verificar sua existência em um banco de dados, assinale a alternativa \textbf{CORRETA}:
    
    \begin{enumerate}[label=\alph*)]
        \item[\textbf{\textcolor{blue}{a)}}] \textcolor{blue}{\textbf{O índice B-Tree organiza os dados de forma ordenada e hierárquica, acelerando muito buscas exatas e por intervalos (como \texttt{BETWEEN}). Podemos verificar os índices de uma tabela usando comandos específicos do SGBD (como \texttt{SHOW INDEX}).}}
        \item O índice B-Tree organiza os dados de forma totalmente aleatória, tornando as consultas (\texttt{SELECT}) mais demoradas, mas acelerando a inserção (\texttt{INSERT}).
        \item Índices B-Tree só podem ser criados em tabelas pequenas. Em tabelas muito grandes, com milhões de registros, o uso de índices é proibido e não é suportado pelo banco de dados.
        \item É impossível listar ou verificar os índices de uma tabela após a sua criação, pois o SGBD esconde essas informações por questões de segurança.
        \item A única maneira de saber se um índice existe é apagando a coluna da tabela e verificando se o banco de dados apresenta um erro de restrição.
    \end{enumerate}

    \vspace{0.5cm}

    \item  Na otimização de bancos de dados, a \textbf{complexidade computacional} (expressa pela notação \textit{Big-O}) é utilizada para descrever o desempenho de algoritmos de busca. Ao procurar um registro específico pelo valor de uma coluna, o SGBD pode realizar uma varredura sequencial completa (\textit{Full Table Scan}) ou consultar um índice baseado em \textbf{Árvore B} (\textit{B-Tree}). Assinale a alternativa que apresenta a complexidade de tempo de busca teórica correta para essas duas operações, respectivamente:
    
    \begin{enumerate}[label=\alph*)]
        \item $O(1)$ para a varredura sequencial e $O(N)$ para a consulta no índice B-Tree.
        \item[\textbf{\textcolor{blue}{b)}}] \textcolor{blue}{\textbf{$O(N)$ para a varredura sequencial e $O(\log N)$ para a consulta no índice B-Tree.}}
        \item $O(\log N)$ para a varredura sequencial e $O(1)$ para a consulta no índice B-Tree.
        \item $O(N^2)$ para a varredura sequencial e $O(N \log N)$ para a consulta no índice B-Tree.
        \item $O(N)$ para ambas as estratégias, pois o índice afeta apenas a leitura física do disco, não alterando a complexidade do algoritmo de busca.
    \end{enumerate}

\end{enumerate}

\end{document}
